\section{Exercise 6}
Implementing the network protocol was particularly difficult for our group, not because it was unusually complex, but because it was something that none of the group members were familiar with.\\
While reading a similar code to gain a better understanding, we quickly noticed that many \texttt{try} and \texttt{catch} blocks were used. Accordingly, we also decided to write catch cases for the methods so that, in the event of an error while communicating with the server, the error would be printed to the console and not just crash.\\
Establishing a connection was relatively simple since setting up such a connection to localhost does not vary much from project to project.\\
Using \texttt{DataInputStream} and \texttt{DataOutputStream} made sending our group ID and reading messages from the server fairly straightforward. The challenge here was getting to know \texttt{DataInputStream} and \texttt{DataOutputStream} and recognizing them as the solution.\\
In the case where we are player 1, we simply return the first valid move. Although there were no difficulties implementing this, as it is not too different from sending our group ID, the fact that we had implemented row and column inversely compared to the server in Exercises 4 and 5 caused some confusion as to why a valid move was repeatedly being rejected by the server. However, this was quickly resolved by simply swapping the two coordinates before sending them to the server. In the future, when we receive moves from the server, we will also swap the coordinates before passing them to our \texttt{GameBoard} file, which is responsible for the game logic.\\
Credits: Fina was responsible for implementing the Network connection in the Java project.